\documentclass[11pt]{article}

\usepackage{amsmath,amssymb}
\usepackage{geometry}
\usepackage{booktabs}

\geometry{margin=1in}

\title{Ramanujan-Type $\pi$ Acceleration and the N-Series Geometric Constant $\pi_g$}
\author{C. Wayne Baker}
\date{June 1, 2026; revised November 2, 2026 (v1.1)}

\begin{document}

\maketitle

\begin{abstract}
Ramanujan's formulas for $\pi$ exhibit extraordinary acceleration through
modular and hypergeometric structure. The N-Series framework approaches the
same constant from a different direction: it begins with constructible
polygonal approximations to the circle and accelerates them by cancelling
successive even-power error terms. This note compares these two viewpoints.
Ramanujan-type formulas reveal a deep analytic acceleration mechanism, while
the N-Series construction defines a geometric limiting constant, denoted
$\pi_g$, through chordal polygonal refinement. The comparison is
complementary rather than competitive: Ramanujan's formulas give exceptional
closed-form series acceleration, while the N-Series gives a constructive
geometric cancellation pathway to $\pi$.
Numerical order tests confirm the expected $2,4,6$ convergence ladder for the
raw, two-scale, and three-scale polygonal approximations.
\end{abstract}

\section{Classical Polygonal Approximation to Pi}

Let
\[
  P_N = N\sin\!\left(\frac{\pi}{N}\right)
\]
denote half the perimeter of the regular $N$-gon inscribed in the unit
circle. Equivalently, $P_N$ is the polygonal chord approximation to the
semicircumference of the unit circle. Since the semicircumference is $\pi$,
one has
\[
  \lim_{N\to\infty} P_N = \pi .
\]

\paragraph{Geometric provenance.}
Although the analytic shorthand $P_N=N\sin(\pi/N)$ contains $\pi$, the
polygonal construction does not use $\pi$ as input. For the doubling scales
used below, $P_N$ is constructed as one half of the perimeter of a regular
inscribed $N$-gon, obtained by repeated chord bisection. The trigonometric
formula is an analytic description of that constructible length.

Expanding the sine function gives
\[
  \sin\!\left(\frac{\pi}{N}\right)
  = \frac{\pi}{N} - \frac{\pi^3}{6N^3} + \frac{\pi^5}{120N^5}
    - \frac{\pi^7}{5040N^7} + \cdots .
\]
Multiplying by $N$, we obtain
\[
  P_N = \pi - \frac{\pi^3}{6N^2} + \frac{\pi^5}{120N^4}
    - \frac{\pi^7}{5040N^6} + \cdots .
\]
Thus the polygonal error has an even-power expansion:
\[
  P_N = \pi + \frac{C_2}{N^2} + \frac{C_4}{N^4} + \frac{C_6}{N^6} + \cdots .
\]
This even-power structure is the basic target of the N-Series cancellation
operators.

\section{The Geometric Constant}

The N-Series viewpoint defines $\pi_g$ as the geometric limit obtained from a
scale-accelerated polygonal sequence. For a fixed integer $k>1$, define
\[
  \widehat{P}^{(k)}_N = P_{kN} + \frac{P_{kN}-P_N}{k^2-1} .
\]
Because the leading polygonal error is proportional to $N^{-2}$, this
two-scale operator cancels the $N^{-2}$ term exactly.

Indeed, suppose
\[
  P_N = \pi + \frac{C_2}{N^2} + \frac{C_4}{N^4} + O(N^{-6}) .
\]
Then
\[
  P_{kN} = \pi + \frac{C_2}{k^2N^2} + \frac{C_4}{k^4N^4} + O(N^{-6}) .
\]
Substituting into the two-scale expression gives
\[
  \widehat{P}^{(k)}_N = \pi + O(N^{-4}) .
\]
The accelerated geometric constant is then
\[
  \pi_g = \lim_{N\to\infty} \widehat{P}^{(k)}_N .
\]
Since $\widehat{P}^{(k)}_N$ tends to $\pi$, it follows that
\[
  \pi_g = \pi .
\]

The significance of $\pi_g$ is not that it defines a different numerical
constant. Rather, it defines $\pi$ through a purely geometric
scale-cancellation process. The constant is recovered from chordal polygonal
geometry, and the acceleration follows from the even-power error structure of
the polygonal approximation.

\section{Higher-Order N-Series Cancellation}

The same idea extends to more scales. Using the three scales $N$, $2N$, and
$4N$, define
\[
  \widehat{P}^{(3)}_N
  = \frac{1}{45}\,P_N - \frac{4}{9}\,P_{2N} + \frac{64}{45}\,P_{4N} .
\]
The weights sum to $1$:
\[
  \frac{1}{45} - \frac{4}{9} + \frac{64}{45} = 1 .
\]
They also cancel the $N^{-2}$ and $N^{-4}$ terms:
\[
  \frac{1}{45} - \frac{4}{9}\cdot\frac{1}{2^2}
  + \frac{64}{45}\cdot\frac{1}{4^2} = 0,
\]
and
\[
  \frac{1}{45} - \frac{4}{9}\cdot\frac{1}{2^4}
  + \frac{64}{45}\cdot\frac{1}{4^4} = 0 .
\]
Consequently,
\[
  \widehat{P}^{(3)}_N = \pi + O(N^{-6}) .
\]

This is the geometric analogue of acceleration. The approximation is not
improved by adding more analytic terms to a series, but by combining the same
geometric construction at multiple scales so that the leading error terms
cancel exactly.

\section{Numerical Validation of the Geometric Acceleration}

To verify the predicted cancellation orders, we computed the raw polygonal
approximation $P_N$, the two-scale approximation $\widehat{P}^{(2)}_N$, and
the three-scale approximation $\widehat{P}^{(3)}_N$ using high-precision
arithmetic. For any approximation $A_N$, the observed order reported at
scale $N$ is
\[
  p_N(A)=\log_2\!\left(
  \frac{|A_{N/2}-\pi|}{|A_N-\pi|}
  \right).
\]
Thus each displayed row compares the errors at scales $N/2$ and $N$. The
observed orders are shown in Table~\ref{tab:pi-g-orders}.

\begin{table}[ht]
\centering
\small
\begin{tabular}{cccc}
\toprule
$N$ & Raw order & Two-scale order & Three-scale order \\
\midrule
$32$   & $1.99791411$ & $3.99875845$ & $5.99923958$ \\
$64$   & $1.99947855$ & $3.99968962$ & $5.99980989$ \\
$128$  & $1.99986964$ & $3.99992240$ & $5.99995247$ \\
$256$  & $1.99996741$ & $3.99998060$ & $5.99998812$ \\
$512$  & $1.99999185$ & $3.99999515$ & $5.99999703$ \\
$1024$ & $1.99999796$ & $3.99999879$ & $5.99999926$ \\
\bottomrule
\end{tabular}
\caption{Observed convergence orders for the raw polygonal approximation, the two-scale N-Series approximation, and the three-scale N-Series approximation.}
\label{tab:pi-g-orders}
\end{table}

The final row gives the observed orders
\[
  1.99999796, \qquad 3.99999879, \qquad 5.99999926,
\]
for the raw, two-scale, and three-scale approximations respectively. This
confirms the predicted cancellation structure:
\[
  P_N-\pi = O(N^{-2}),
\]
\[
  \widehat{P}^{(2)}_N-\pi = O(N^{-4}),
\]
and
\[
  \widehat{P}^{(3)}_N-\pi = O(N^{-6}) .
\]

For completeness, the corresponding absolute errors at $N=1024$ are
\[
  4.92831263353783379\times 10^{-6},
\]
\[
  5.79840008299119192\times 10^{-13},
\]
and
\[
  8.12154741379365833\times 10^{-21} .
\]
Thus the numerical data support the theoretical interpretation of $\pi_g$:
the geometric constant is recovered not merely by polygonal convergence, but
by a scale-cancellation ladder that removes successive even-power error terms.

\section{Ramanujan-Type Pi Acceleration}

Ramanujan's formulas for $\pi$ arise from a very different source. His series
for $1/\pi$ are connected to modular equations, elliptic integrals, and
hypergeometric functions. A representative Ramanujan-type structure is
\[
  \frac{1}{\pi} = \sum_{n=0}^{\infty} A_n,
\]
where the coefficients $A_n$ contain factorial or hypergeometric factors
together with rapidly decaying modular powers.

A famous example from Ramanujan's 1914 work has the form
\[
  \frac{1}{\pi}
  = \frac{2\sqrt{2}}{9801}
    \sum_{n=0}^{\infty}
    \frac{(4n)!\,(1103+26390n)}{(n!)^4\,396^{4n}} .
\]
This formula gains many digits per term. Its acceleration is not polygonal or
chordal; it comes from the analytic and modular structure behind the series.
Ramanujan acceleration is modular and analytic, whereas N-Series acceleration
is geometric and scale-cancellative.

section{Quantitative comparison and claim boundary}
For Ramanujan's displayed series, write
\[
t_n=\frac{(4n)!(1103+26390n)}{(n!)^4\,396^{4n}}.
\]
Then
\[
\frac{t_{n+1}}{t_n}=
\frac{(4n+1)(4n+2)(4n+3)(4n+4)}{(n+1)^4\,396^4}
\frac{1103+26390(n+1)}{1103+26390n},
\]
so the limiting ratio is \(1/99^4\), corresponding to
\(4\log_{10}(99)\approx7.98\) decimal digits per summand.
This term-wise geometric reduction differs from the N-Series
algebraic orders 2, 4 and 6 as polygonal resolution increases.
The two are not comparable cost metrics.

On the doubling ladder \(N=4\cdot2^j\), the polygon lengths
and their finite rational linear combinations are constructible.
The limit \(\pi_g=\pi\) is not thereby constructible.
The polygonal expansion, Richardson cancellation, and Ramanujan
identity are classical. No priority, new modular identity,
or performance superiority is claimed here.

\section{Comparison of the Two Viewpoints}

The two approaches should not be treated as competing approximations in the
ordinary numerical sense. Ramanujan's formula is vastly faster than raw
polygonal approximation and remains one of the most remarkable analytic
accelerations for $\pi$. The N-Series contribution is different: it shows that
high-order acceleration can also arise from exact cancellation inside a purely
geometric sequence.

The comparison is therefore structural. Ramanujan's method reveals that $\pi$
is accessible through modular and hypergeometric identities. The N-Series
method reveals that $\pi$ is also accessible through constructible polygonal
geometry equipped with scale-cancellation operators.

In the Ramanujan setting, the acceleration is hidden in the modular
transformation. In the N-Series setting, the acceleration is visible in the
moment conditions imposed on the scale weights. For the three-scale operator,
these conditions are
\[
  \sum_{j=0}^{2} w_j = 1,
  \qquad
  \sum_{j=0}^{2} \frac{w_j}{m_j^{2}} = 0,
  \qquad
  \sum_{j=0}^{2} \frac{w_j}{m_j^{4}} = 0,
\]
with scale factors
\[
  m_0=1, \qquad m_1=2, \qquad m_2=4 .
\]
This cancellation structure is algebraic, but it acts on a geometric sequence.

\section{Interpretation}

The N-Series does not improve Ramanujan's formula in the sense of producing a
faster numerical series for $\pi$. Ramanujan's formulas are already
extraordinarily efficient. Instead, the N-Series provides a complementary
geometric interpretation of acceleration: the leading error terms in a
constructible sequence can be removed exactly by combining multiple scales.

This distinction is important for positioning the work. The claim is not that
polygonal acceleration replaces modular acceleration. The claim is that both
expose hidden structure. Ramanujan exposes hidden analytic structure in
elliptic and modular functions; the N-Series exposes hidden even-power error
structure in geometric approximations.

Thus the natural bridge is
\[
  \text{Ramanujan acceleration}
  \quad\longleftrightarrow\quad
  \text{N-Series acceleration},
\]
where Ramanujan's acceleration is analytic and modular, while the N-Series
acceleration is geometric and constructive.
\section{Limitations}

This note is a structural comparison, not a claim that polygonal N-Series
acceleration competes numerically with Ramanujan-type series.  Ramanujan's
formulas are analytic, modular, and extraordinarily fast.  The N-Series
construction has a different purpose: it shows that exact high-order
acceleration can also arise from geometric scale cancellation.

The numerical tests verify the predicted cancellation orders for the tested
polygonal sequence, but they do not replace the analytic theory behind
Ramanujan's formulas.  The comparison is therefore complementary: one
mechanism is modular and hypergeometric, while the other is geometric and
scale-cancellative.

\section{Conclusion}

Ramanujan-type formulas and the N-Series construction give two distinct
pathways to the same constant $\pi$. Ramanujan's pathway is analytic, modular,
and hypergeometric. The N-Series pathway is geometric, chordal, and
constructive. In the N-Series setting, $\pi_g$ is not a new constant but a
geometric recovery of $\pi$ through scale-cancelled polygonal approximations.

The significance of this comparison is complementary. Ramanujan's formulas
demonstrate the power of modular analytic acceleration, while the N-Series
demonstrates that exact multi-scale cancellation can also accelerate purely
geometric constructions. Together, they point to a broader principle: rapid
convergence often arises when a hidden structure allows the leading error
terms to be cancelled rather than merely reduced.

\begin{thebibliography}{9}

\bibitem{ramanujan}
S. Ramanujan,
``Modular equations and approximations to $\pi$,''
\emph{Quarterly Journal of Mathematics}, 45, 1914, pp.\ 350--372.

\bibitem{borwein}
J. M. Borwein and P. B. Borwein,
\emph{Pi and the AGM: A Study in Analytic Number Theory and Computational
Complexity}.
Wiley, 1987.

\bibitem{richardson}
L. F. Richardson,
``The approximate arithmetical solution by finite differences of physical
problems,''
\emph{Philosophical Transactions of the Royal Society A}, 210, 1911,
pp.\ 307--357.

\bibitem{archimedes}
Archimedes,
\emph{Measurement of a Circle}.

\end{thebibliography}

\end{document}
